\section{Introduction}
\label{sec:intro}

Peripheral nerve stimulation (PNS) underpins clinical therapies for a
growing range of indications including drug-resistant epilepsy,
treatment-resistant depression, rheumatoid arthritis, heart failure
and chronic pain~\cite{howland2014vns,koopman2016vns,bonaz2017vagus,deferrari2011cardiofit}.
Because peripheral nerves carry mixed populations of motor, sensory
and autonomic fibers organised into anatomically distinct fascicles
whose number, calibre and somatotopic ordering vary substantially
across species and individuals~\cite{settell2020vagotopy,pelot2023humanvagus},
the therapeutic question is rarely \emph{whether} to stimulate but
\emph{which subset} of fibers to engage---a problem of
\emph{selectivity}~\cite{aristovich2021,thompson2023spatially,musselman2023}.  Two design
levers drive selectivity: (i) the spatial distribution of cathodic
and anodic charge across multi-contact cuff electrodes, and (ii) the
temporal shape of the stimulus waveform, both of which couple
nonlinearly to fiber threshold via the extracellular
field and the active membrane dynamics of each
fiber~\cite{rattay1989analysis,grill1995stimulus,wongsarnpigoon2010energy}.

Computational evaluation of candidate stimulation patterns has
become a standard component of PNS device development.  The dominant
workflow links a finite-element solve for the extracellular potential
$V_e(\mathbf{x})$ in a realistic nerve geometry to compartmental
NEURON~\cite{hines1997neuron} simulations of canonical axon models---
typically the McIntyre--Richardson--Grill (MRG) double-cable model (or
the simpler Sweeney model) for myelinated
fibers~\cite{mcintyre2002,sweeney1987} and the Sundt model for
unmyelinated C-fibers~\cite{sundt2015}---with
parameter scanning to map activation thresholds across fascicles.
Open-source pipelines such as ASCENT~\cite{musselman2021ascent} and
pyfibers~\cite{hussain2024} have made this workflow accessible, but
they share a fundamental cost: every candidate design requires
\emph{thousands} of independent NEURON forward simulations, run
serially over CPU cores, with no mechanism to compute the gradient of
a selectivity objective with respect to design variables.  Optimisation
over the resulting cost landscape is consequently restricted to
gradient-free methods---grid sweeps, finite-difference
approximations~\cite{hussain2024} or evolutionary
search~\cite{aristovich2021}---scaling poorly to the high-dimensional
spaces opened up by waveform shaping or many-contact cuffs.
A second, less-examined economy compounds the first: to keep the fibre
count manageable, these pipelines typically reduce each fascicle to a
single representative (centroid) fibre~\cite{hussain2024}.  Whether
this reduction preserves the selectivity of candidate designs --- or
quietly biases it --- has not been tested, because testing it demands
simulating the full fibre population under the very fields the reduced
model selects.

Recent work in differentiable simulation suggests an alternative.
Frameworks such as \texttt{jaxley}~\cite{deistler2024} have shown that
the Hodgkin--Huxley membrane equations and multi-compartment cable
dynamics can be expressed in JAX~\cite{jax2018} so that gradients
of arbitrary scalar functions of the membrane trajectory are
available via reverse-mode automatic differentiation, all within a
single GPU-vectorised forward pass.  In the central nervous system
context, this has enabled gradient-based inference of channel
densities and morphological parameters at unprecedented scale.
For peripheral nerve fiber models, however, the differentiable
implementation does not yet exist: the MRG double-cable model
requires a coupled $(V_i, V_{\mathrm{px}})$ backward-Euler integration
that is not provided out-of-the-box, the other models (the myelinated
Sweeney and the unmyelinated Sundt and Rattay) each carry
calibration-sensitive constants that must match NEURON's analytic
solver to bench-marking accuracy~\cite{musselman2023}, and selectivity
optimisation requires a batched-fiber forward solver that vectorises
across large fiber populations on a single device.  Building the differentiable peripheral nerve
forward solver therefore reduces to three engineering questions:
(i) can the coupled MRG cable equations be expressed in JAX with
sufficient numerical accuracy that the resulting simulator passes the
same NEURON-comparison validation suite that ASCENT and pyfibers use;
(ii) is the resulting solver competitive with NEURON on throughput
when batched across large fiber populations; and (iii) does the
differentiability enable selectivity optimisation strategies that
the gradient-free pipelines cannot perform?

In this paper we use population-scale differentiable simulation to ask
whether the reduced-order fibre sampling that keeps selectivity models
tractable biases the selectivity they predict --- and we find that it
does, systematically and in a species-dependent way.  Our
contributions are:

\begin{enumerate}
  \item \textbf{The central finding.}  On a cross-species cohort of
        microCT-segmented vagus nerves, stimulation configurations
        optimised on sparsely-sampled fibres --- the per-fascicle
        reduction used in GPU-accelerated selectivity
        work~\cite{hussain2024} --- lose substantial selectivity when
        deployed on the full fibre population.  This \emph{deployment
        penalty} scales with the number of fibres per target fascicle
        and thereby differs systematically between species (small in
        swine, large in human), exposing a translational bias that
        reduced-order models cannot detect
        (Section~\ref{sec:results-sparse}).
  \item \textbf{The enabling simulator.}  \emph{jaxon}, a fully
        differentiable, GPU-native reimplementation in JAX of four
        canonical axon models (myelinated MRG double-cable and Sweeney;
        unmyelinated Sundt and Rattay), validated against NEURON on
        strength--duration, conduction-velocity and intracellular-trace
        benchmarks ($99.6\,\%$ of 943 configurations within $1\,\%$)
        and $\sim\!820\times$ faster than \texttt{pyfibers} at
        $N=100{,}000$ fibres on a single A100 --- fast enough to
        optimise and re-score whole-nerve populations
        (Sections~\ref{sec:results-validation}--\ref{sec:results-scaling}).
  \item \textbf{The pipeline.}  An end-to-end differentiable
        selectivity-optimisation method over per-contact stimulus
        amplitudes, exposing exact gradients that gradient-free PNS
        pipelines lack, demonstrated on 29 microCT-derived vagus
        anatomies with FEM-computed extracellular fields --- to our
        knowledge the first cross-species gradient-based selectivity
        optimisation on real microCT nerve geometry
        (Section~\ref{sec:results-duke}).
\end{enumerate}

We discuss the limits of the current implementation, including the
anatomical-ceiling cases encountered on the Duke microCT cohort where
the cuff cannot reach the assigned target fascicle at any amplitude
(Section~\ref{sec:discussion}).
